\emph{This chapter was written by Claude, Anthropic's AI, from an
outline and direction by Carsten Wulff, who reviewed and edited the
result. The figures are Claude's, in the book's style. The commit
history of the book's repository records precisely who wrote what.}

Circuit theory is Maxwell's equations with the fine print deleted. The
deletion is a good deal - we get Kirchhoff, nodes and branches, and we
can design a chip without solving a single boundary value problem. But
the fine print does not go away, and every so often it sends a bill: a
supply that bounces, a clock that couples into an ADC, an inductor that
is mostly resistor, a wire that has become an antenna.

This chapter restores exactly three pieces of the fine print, the three
an IC designer keeps paying for: currents close in loops, every gap is a
capacitor, and every loop is an inductor. Radiation - the part everyone
treats as magic - is just what the fine print does at high frequency.

There will be no vector calculus gymnastics here. We use the integral
forms, in words, and tie every claim to a chapter of this book where you
will meet it again.

\section{The four equations, in
words}\label{the-four-equations-in-words}

\[ \oint \vec{E} \cdot d\vec{A} = \frac{Q}{\varepsilon_0} \]

\[ \oint \vec{B} \cdot d\vec{A} = 0 \]

\[ \oint \vec{E} \cdot d\vec{\ell} = -\frac{d\Phi_B}{dt} \]

\[ \oint \vec{B} \cdot d\vec{\ell} = \mu_0 I + \mu_0 \varepsilon_0 \frac{d\Phi_E}{dt} \]

Gauss: electric field lines start on positive charge and end on negative
charge. Count the lines leaving a closed surface and you have counted
the charge inside.

Gauss for magnetism: magnetic field lines do not start or end anywhere.
There is no magnetic charge; every field line is a closed loop.

Faraday: a changing magnetic flux through a loop drives a voltage around
it. This is the only way to make an electric field whose lines close on
themselves rather than ending on charge.

Ampere, with Maxwell's correction: currents make magnetic fields curl
around them - and so does a \emph{changing electric field}. The second
term is the fine print that makes capacitors, and radio, work. It is
called the displacement current, and it deserves a better reputation.

One more law hides inside these four. Take Ampere's equation and close
the surface: the conduction current in, plus the displacement current
in, must equal zero. Charge is conserved, always, everywhere. That
innocent bookkeeping statement is the most useful sentence in this
chapter.

\section{Currents always run in
loops}\label{currents-always-run-in-loops}

Charge conservation says the total current into any closed surface is
zero. Draw a surface around any point of your circuit: whatever current
comes in must leave. Follow it, and you must eventually come back to
where you started. \textbf{Every current closes a loop.} There are no
exceptions - not for signals, not for supplies, not for that
``unidirectional'' clock trace.

The design consequence is that the return path is not optional. You only
get to choose \emph{where} it goes. Draw the loop deliberately, as in
Figure 1(a), and you know its area, its inductance and its victims.
Forget it, as in Figure 1(b), and the current still closes - through the
substrate, through a neighbouring supply, through whatever shared ground
is available - and everything on that accidental path sees your signal
as ground bounce.


{
\centering
\aicfig{media/mx_loop_tikz.png}
}

\emph{Figure 1: Every current closes a loop. (a) The loop you drew. (b)
The loop you got: the forgotten return closes through the shared ground,
and everything on that path sees it}

When you debug a noisy chip, do not ask ``where does this signal go?''.
Ask ``where does it come back?''. The second question finds the problem.

\textbf{Don't ask ``where does the signal go?''}

\textbf{Ask ``where does it come back?''}

\section{The capacitor falls out}\label{the-capacitor-falls-out}\index{the capacitor falls out@The capacitor falls out}

Now break the top wire of the loop, as in Figure 2. Kirchhoff panics -
the circuit is open. Maxwell does not: charge piles up on the two faces
of the break, an electric field grows between them, and the changing
field \emph{is} a current. Ampere's correction term
\(\partial D/\partial t\) carries the loop across the gap as if the wire
were never cut. Charge conservation is satisfied at every instant; the
loop never broke.

A capacitor is a deliberately good gap: two large faces, close together,
with a dielectric that multiplies the effect. The \(I = C\, dV/dt\) you
have used since your first circuits course is Ampere's displacement
current wearing a component symbol.

And because \emph{any} gap does this, every pair of conductors on your
chip is a capacitor whether you asked for one or not - the parasitics
chapter is the bill.


{
\centering
\aicfig{media/mx_cap_tikz.png}
}

\emph{Figure 2: Break the conductor and the loop continues through the
gap as displacement current. A capacitor is a deliberately good gap}

\section{The inductor falls out}\label{the-inductor-falls-out}\index{the inductor falls out@The inductor falls out}

Figure 3(a) walks the three steps. \textbf{One:} push a rising current
through a loop of wire, and Ampere says that current wraps a magnetic
field around the wire. \textbf{Two:} some of that field threads the loop
itself - the flux \(\Phi\) is proportional to the current, with a
constant decided purely by the geometry, and we call it \(L\).
\textbf{Three:} Faraday says a changing flux drives a voltage around the
loop, and Lenz's rule says it drives it in the direction that opposes
the change. The loop fights you, and at the terminals that reads
\(v = L\, di/dt\): ramp the current, and the loop answers with a
voltage, exactly as the inset shows.

An inductor is a deliberately good loop. Wind the same loop \(N\) times,
Figure 3(b), and the bargain improves twice over: \(N\) turns make \(N\)
times the flux, and every turn feels all of it - which is why the
inductance grows as \(N^2\).

The converse of the capacitor's lesson holds here: every loop on your
chip is an inductor whether you asked or not. Since every current runs
in a loop, every current has an inductance. That is why fast edges ring
- the loop's \(L\) and the gap's \(C\) are always both present, and
together they make a resonator you never drew.


{
\centering
\aicfig{media/mx_ind_tikz.png}
}

\emph{Figure 3: Every loop encloses flux, and a changing current fights
its own flux. An inductor is a deliberately good loop; a coil is the
same loop N times}

\section{Every wire is both}\label{every-wire-is-both}\index{every wire is both@Every wire is both}

So the fine print reads: every gap is a capacitor, every loop is an
inductor, and every current must loop. Put the three together over a
ground plane and you can predict something that surprises most people
the first time: \emph{where the return current flows}.

At DC the return current spreads across the whole plane - it minimizes
resistance, Figure 4(a). At high frequency it crowds into a narrow band
directly under the signal wire, Figure 4(b). It is minimizing impedance,
and above a few megahertz the impedance is dominated by the loop
inductance - which shrinks with the loop area. The current chooses the
smallest loop available.

This one picture is most of signal integrity. Slot the ground plane
under a trace and the return detours around the slot: the loop area
explodes, the inductance with it. Decoupling works the same way: the
capacitor's value matters less than the loop area between it and the
load, because it is the loop inductance that decides how fast the
capacitor can deliver charge.


{
\centering
\aicfig{media/mx_return_tikz.png}
}

\emph{Figure 4: Where the return current flows in a ground plane. (a) At
DC it spreads to minimize resistance. (b) At high frequency it hugs the
signal wire to minimize loop inductance}

\section{The antenna does not
radiate}\label{the-antenna-does-not-radiate}

Here is the claim this chapter has been building towards: the metal of
an antenna does not radiate. The metal only sets up boundary conditions.
The drive sloshes charge up and down the arms; the charge makes fields
around the structure; and it is the \emph{fields} that carry the power
away.

Close to the dipole, Figure 5(a), the electric field lines run from one
arm to the other. Every half cycle the drive reverses, the field lines
collapse back onto the metal, and their energy returns to the circuit.
This is the near field: energy borrowed and repaid, which at the
terminals looks like reactance - the antenna below resonance is just a
capacitor.

But the news that the drive has reversed travels at the speed of light.
Field lines further than about \(\lambda/2\pi\) from the arms get the
news too late: the charge that anchored them has already moved on, and
the lines have nothing to end on. Gauss offers them one way out - close
on yourself. Figure 5(b): the loops detach, each one a ring of
displacement current sustaining the magnetic field of the next (Ampere),
which sustains the electric field of the next (Faraday), and the pair
leapfrogs outward at \(c\) with no metal anywhere. That self-sustaining
leapfrog \emph{is} the electromagnetic wave.

At the antenna terminals, the energy that left with the detached loops
never comes back. Power delivered and not returned looks, to the
circuit, exactly like a resistor: the radiation resistance. It is the
receipt for the escaped loops.


{
\centering
\aicfig{media/mx_antenna_tikz.png}
}

\emph{Figure 5: (a) Near the dipole the field lines end on the metal and
their energy returns every half cycle. (b) Beyond about lambda over 2 pi
the lines close on themselves - displacement current with no metal - and
carry power away at c}

The wavelength sets the geometry. Figure 6 shows the workhorse as it
actually ships: a quarter-wave copper trace on a PCB, fed by the radio
chip, over a keep-out where the ground is cut away. The ground pour is
the other half of the antenna - the trace's image in it completes the
dipole. The current standing wave is maximum at the feed and zero at the
tip, and resonance lands at \(l = \lambda/4\): at 2.4 GHz about 31 mm,
which is why the antenna region of a Bluetooth board is the size it is,
and why nothing on a millimetre scale chip radiates \emph{efficiently}
by accident. Accidental radiators are inefficient antennas - but a
receiver channel fighting for -100 dBm does not need your clock harmonic
to be efficient, only present.


{
\centering
\aicfig{media/mx_monopole_tikz.png}
}

\emph{Figure 6: The quarter-wave monopole as it ships, seen from above:
the antenna trace on its keep-out, fed by the radio, with the ground
pour as the other half of the antenna. Above, on the same scale, the
current standing wave that sets l = lambda/4 - about 31 mm at 2.4 GHz}

\section{What this buys you on-chip}\label{what-this-buys-you-on-chip}\index{what this buys you on-chip@What this buys you on-chip}

Everything in this chapter reappears later in the book wearing a
different costume.

On-chip inductors are poor because the fine print is against them twice:
the loops that fit on a die are small, so \(L\) is small, and they sit
on a conductive substrate, so Faraday drives eddy-current losses in
exactly the silicon we paid so much for. The LC oscillator chapter lives
with the resulting Q.

Decoupling capacitors are loop design, not capacitor selection: the
inductance of the loop from capacitor to load sets the frequency above
which the capacitor stops helping.

Supply and ground bounce are Figure 1(b) at chip scale, and the cure is
the same as the diagnosis: give every fast loop a small, deliberate
return.

And the radio chapter's antenna, matching network and Friis budget are
this chapter run in reverse: arrange the boundary conditions so the
detaching loops are not an accident but the product.

\begin{itemize}
\tightlist
\item
  On-chip inductors: small loops, lossy substrate - the fine print
  charges twice
\item
  Decoupling is loop design, not capacitor selection
\item
  Ground bounce is Figure 1(b) at chip scale
\item
  A radio is this chapter, run on purpose
\end{itemize}

\section{Summary}\label{summary}

The one-page version of this chapter:

\begin{itemize}
\tightlist
\item
  Circuit theory is Maxwell with the fine print deleted; the fine print
  still bills you
\item
  Charge conservation means every current closes a loop - the return
  path is not optional, only its location is
\item
  A capacitor is a deliberately good gap: displacement current carries
  the loop across
\item
  An inductor is a deliberately good loop: every loop encloses flux, and
  a changing current fights its own flux
\item
  At high frequency the return current hugs the signal wire, because the
  smallest loop has the smallest inductance
\item
  The antenna metal only sets boundary conditions: field loops that
  detach beyond lambda/2pi carry the power, and radiation resistance is
  their receipt
\end{itemize}

\section{Would you like to know
more?}\label{would-you-like-to-know-more}

The best intuition-first treatment of these ideas remains Feynman's
Lectures on Physics, Volume II - chapter 18 for the full set of
equations and what they mean, and chapter 24 onward for waveguides and
radiation. For the signal-integrity consequences, Howard Johnson's
High-Speed Digital Design (the ``black magic'' book) is the field manual
for Figure 4.

\begin{itemize}
\tightlist
\item
  Feynman Lectures on Physics, Volume II,
  \href{https://www.feynmanlectures.caltech.edu/II_18.html}{chapter 18}
\item
  Howard Johnson \& Martin Graham, High-Speed Digital Design: A Handbook
  of Black Magic
\end{itemize}
